\section{Round-twenty-one positive closure: fixed-size smoothing blocks and exact-number saddle extraction}
\label{sec:r21-b1}

This section replaces the former Round-Seventeen source.  The scaling is the
three-dimensional Boltzmann--Grad regime of B2:
$\mu_\varepsilon=\varepsilon^{-2}$ and
$N/\mu_\varepsilon\to\rho$ in a fixed compact subset of $(0,\infty)$.  The
finite-volume law is the positive exact-$N$ hard-sphere phase-space integral.
Connected activities are used only as analytic estimates for this positive
integral and are never interpreted as independent random components.

\subsection{Constraint space and the source-dependent saddle}

Let
\[
 C(x,v)=\bigl(v,|v|^2/2,\chi_1(x),\ldots,\chi_m(x)\bigr),
\]
where the smooth cell functions $\chi_j$ have compactly supported gradients in
pairwise separated interior boxes.  Let $\mathcal R$ be the finite-dimensional
space of affine relations among the components of $C$ and put
$E=\mathbb R^d/\mathcal R$.  All covariances and Fourier variables below are
on $E$; redundant partition-of-unity coordinates are not assigned a positive
variance.

For a bounded regular path/contact source $H$ and multiplier
$\lambda\in E^*$, define
\[
 Z_{\varepsilon,N}(H,\lambda)=
 {1\over N!}\int_{\mathcal D_{\varepsilon,N}\times\mathbb R^{3N}}
 \exp\left\{\sum_{i=1}^N\lambda\cdot C(x_i,v_i)
              +\mu_\varepsilon H(\mathbf X)\right\}
 d\mathbf x\,d\mathbf v .
\]
The source class is the same analytic ball used in the B2 pressure theorem.

\begin{theorem}[Uniform canonical pressure and saddle]
\label{thm:r21-b1-saddle}
There is an open convex multiplier domain $\mathcal O\subset E^*$ such that
\[
 q_N(H,\lambda)={1\over N}\log Z_{\varepsilon,N}(H,\lambda)
\]
converges, with three multiplier derivatives and two source derivatives,
locally uniformly to a strictly convex function $q(H,\lambda)$.  On every
compact set $A$ in the interior of its mean image, the equation
\[
 D_\lambda q(H,\lambda)=a
\]
has a unique solution $\lambda(H,a)$, analytic in $H$ and $a$, and
\[
 cI_E\le D_\lambda^2q(H,\lambda(H,a))\le CI_E.
\]
The same estimates hold for $q_N$ and its saddle for all sufficiently small
$\varepsilon$.
\end{theorem}

\begin{proof}
Insert an activity variable $z$ in the B2 grand-canonical pressure and use its
normally convergent connected series on a short time slice.  The coefficient
of $z^N$ is exactly the positive integral above.  Uniform convergence on a
complex annulus gives the canonical logarithmic limit and permits the stated
source and multiplier derivatives.  At zero interaction, the log moment of
$C$ is strictly convex on the quotient $E$ because no nonzero element of
$E^*$ is constant on the support.  The B2 connected remainder has Hessian
$O(\varepsilon)$ uniformly on compact multiplier sets; hence the lower and
upper covariance bounds persist.

For an exponential family on a compact multiplier chart, the gradient of a
strictly convex log partition is a diffeomorphism onto the interior of the
convex support.  Properness follows because approaching the chart boundary
concentrates on an exposed face, while $A$ stays a positive distance from all
faces.  The interacting gradient is a uniform small perturbation and has the
same degree.  Thus it is globally one-to-one and onto $A$, not merely locally
invertible.  The Banach analytic implicit-function theorem gives dependence
on $H$ and $a$.
\end{proof}

\subsection{Fixed-size full-rank smoothing blocks}

Choose $b=d+2$ pairwise separated position boxes and velocity balls so that
the block map
\[
 \mathcal C_B:(x_1,v_1,\ldots,x_b,v_b)\longmapsto
 \sum_{j=1}^b C(x_j,v_j)\in E
\]
has a $d$-dimensional minor bounded below on a compact subbox.  Position
variables are included in this minor whenever the Fourier direction is
position dominated.  The boxes are separated by a fixed distance, so no
hard-core boundary occurs inside the block chart.

A block is called fresh if its particles have no prescribed collision with
previously exposed particles during the short preparation layer and their
initial data lie in one of the full-rank charts.  Expose blocks successively,
conditioning on the exterior configuration.

\begin{lemma}[Uniform conditional block regularity]
\label{lem:r21-b1-block}
There are an integer $s>d+5$, constants $c,C>0$, and $\alpha>0$ such that the
following holds.  Except on an event of probability at most $e^{-cN}$, at
least $\alpha N$ disjoint fresh blocks can be selected.  Conditional on the
exterior and on all previously selected blocks, the next block has density
\[
 p_B=p_B^{(0)}R_B
\]
on a full-rank chart, where $p_B^{(0)}$ is a product density bounded above and
below and
\[
 \sum_{|\gamma|\le s}\|\partial_B^\gamma R_B\|_{L^1}\le C.
\]
The constants are uniform in $N,\varepsilon,H$, and the saddle.
\end{lemma}

\begin{proof}
Tile phase space by a fixed finite family of separated position--velocity
boxes.  Hard-core exclusion removes $O(N\varepsilon^3)=O(\varepsilon)$ of the
conditional position volume because $N\asymp\varepsilon^{-2}$.  The B2
connected expansion with one block of marked root variables gives, after
conditioning on the exterior, a Radon--Nikodym factor whose derivatives in
that fixed block are bounded by the normally convergent marked cluster
series.  The block size $b$ and derivative order $s$ are fixed; no derivative
order grows with $N$.  The dynamic source is differentiated only through
clusters meeting the block, and the B2 source ball gives the same bound.

On each exposure, either a full-rank box is available with conditional
probability at least $p_*>0$, or the configuration lies in a union of
forbidden-volume and singular-coordinate events of conditional probability
$O(\varepsilon)$.  A martingale Chernoff bound for these bounded conditional
success indicators gives at least $\alpha N$ successes with probability
$1-e^{-cN}$.  Taking the finite union over the block atlas proves the lemma.
\end{proof}

\begin{lemma}[One-block Fourier smoothing]
\label{lem:r21-b1-oneblock}
For every conditional block density in
Lemma~\ref{lem:r21-b1-block},
\[
 \left|\int e^{iu\cdot\mathcal C_B}p_B\right|
 \le C(1+|u|)^{-s}
\]
for $|u|\ge1$, uniformly in the direction $u\in E^*$.  On every compact
annulus $\delta\le|u|\le M$ the left side is at most $1-c_{\delta,M}$.
\end{lemma}

\begin{proof}
Partition the unit sphere in $E^*$ according to a nonvanishing minor of
$D\mathcal C_B$.  On each piece use the corresponding $d$ block coordinates
as coarea coordinates.  The push-forward density is $W^{s,1}$ with a uniform
norm by Lemma~\ref{lem:r21-b1-block}; $s$ integrations by parts in the
push-forward variable give the first estimate.  This integration includes
position coordinates in position-dominated directions, so velocity
integration is never asked to detect a position-only phase.  Compactness and
full-dimensional support imply strict modulus below one on the annulus; the
uniform density bounds make the gap uniform.
\end{proof}

\subsection{A global characteristic majorant without $O(N)$ derivatives}

Let $\Phi_{N,H,a}(u)$ be the centered characteristic function under the exact
canonical saddle law.  Successively integrate the fresh blocks selected in
Lemma~\ref{lem:r21-b1-block}.  At each step only derivatives in the current
fixed-size block are used.

\begin{theorem}[Blockwise global Fourier estimate]
\label{thm:r21-b1-fourier}
Uniformly in the source and target charts, there are $c,C,\delta>0$ such that
\[
\begin{array}{ll}
|u|\le\delta:&
 |\Phi_N(u)|\le Ce^{-cN|u|^2},\\[1mm]
\delta<|u|\le e^{cN}:&
 |\Phi_N(u)|\le Ce^{-cN},\\[1mm]
|u|>e^{cN}:&
 |\Phi_N(u)|\le C(1+|u|)^{-s\alpha N/2}.
\end{array}
\]
In particular, for one fixed $s_0>d+4$,
\[
 |\Phi_N(u)|\le C e^{-cN}(1+|u|)^{-s_0}
 +C(1+|u|)^{-s_0-cN}
\]
outside a neighborhood of zero.  All derivatives used in the proof have
order at most the fixed number $s$ in each block.
\end{theorem}

\begin{proof}
The small-frequency line follows from the uniform covariance lower bound and
a fourth-cumulant estimate obtained by differentiating the B2 connected
pressure.  On a compact annulus, condition and integrate one fresh block at a
time.  Lemma~\ref{lem:r21-b1-oneblock} gives a factor at most $1-c$; iterating
through $\alpha N$ blocks gives $e^{-cN}$.  For large $|u|$, the same sequential
Fubini argument uses the polynomial factor $(1+|u|)^{-s}$ from each block.
The derivative is always taken only in the current block conditional density;
no mixed derivative across all blocks is generated.  The exceptional event is
covered by the remaining block charts rather than bounded by a frequency
independent number: the measurable greedy selection partitions the entire
regular configuration space into finitely many chart families, each with at
least $\alpha N/2$ full-rank blocks.  The singular chart boundaries have zero
canonical measure and are absorbed by a smooth partition of unity.  Combining
the finite atlas gives the final two-line majorant.
\end{proof}

\subsection{Mixed lattice--continuous coefficient asymptotics}

Separate the genuinely lattice coordinates of the constraint from the
continuous quotient coordinates, writing $E=E_{\mathbb Z}\oplus E_{\mathbb R}$.
For a target $a$ in the compact interior mean image, let
$\Sigma_{H,a}$ be the covariance at the source-dependent saddle.  An
admissible shell is
\[
 S_N(k,D)=\{C_N^{\mathbb Z}=k,
 \ N^{-1/2}(C_N^{\mathbb R}-Na_{\mathbb R})\in D\},
\]
where $D$ has uniformly $C^2$ boundary and positive Gaussian mass.  Shrinking
families are allowed when their Gaussian mass dominates the uniform inversion
error.  No exact point event is asserted for a continuous coordinate.

\begin{theorem}[Uniform local coefficient theorem]
\label{thm:r21-b1-coefficient}
Uniformly in $H$, $a$, admissible lattice targets, and admissible shells,
\[
 \mathbb P_{H,\lambda(H,a)}(S_N(k,D))
 ={e^{-\frac1{2N}z_{\mathbb Z}^T
       \Sigma_{\mathbb Z}^{-1}z_{\mathbb Z}}
   \over(2\pi N)^{d_{\mathbb Z}/2}
       \sqrt{\det\Sigma_{\mathbb Z}}}
 \left[\gamma_{\Sigma_{\mathbb R}}(D)+o(1)\right].
\]
Before tilting, the right side is multiplied by the exact Cram\'er factor
$e^{-NI_H(a)}$.  The theorem remains valid with a bounded regular
path/contact insertion and with two source derivatives.
\end{theorem}

\begin{proof}
Use Fourier inversion on the torus for the lattice coordinates and on
$E_{\mathbb R}^*$ for the continuous coordinates.  The small ball gives the
Gaussian by the cubic pressure expansion.  The complement is integrable and
$o(N^{-d/2})$ by Theorem~\ref{thm:r21-b1-fourier}.  Approximate
$\mathbf1_D$ from inside and outside by smooth functions supported in a
boundary layer of width $h_N$.  Gaussian perimeter bounds make the layer mass
$O(h_N)$; choose $h_N$ larger than the inversion error and smaller than the
shell inradius.  Relative error is claimed only when the Gaussian shell mass
is larger than this absolute error.  A path/contact insertion is a bounded
source derivative of the same characteristic function, so the estimates are
unchanged.
\end{proof}

\subsection{Exact-number extraction from the B2 pressure}

Let $\mathcal Q(H,z,\lambda)$ be the B2 grand-canonical pressure with activity
$z$.  At the desired density and constraint, solve jointly
\[
 \partial_{\log z}\mathcal Q=\rho,
 \qquad D_\lambda\mathcal Q=a.
\]
Strict convexity on $\mathbb R\oplus E^*$ gives the unique saddle
$(z(H,a),\lambda(H,a))$.

\begin{theorem}[Source-uniform exact-number pressure]
\label{thm:r21-b1-extraction}
The exact coefficient satisfies
\[
 [z^N]Z_\varepsilon(H,z,\lambda)
 =\exp\{\mu_\varepsilon\mathcal Q(H,z_*,\lambda_*)
       -N\log z_* -N I_H(a)+o(\mu_\varepsilon)\}
 \mathcal G_{\varepsilon,N}(H,a),
\]
where $\mathcal G_{\varepsilon,N}$ has the Gaussian lattice--continuous
asymptotic of Theorem~\ref{thm:r21-b1-coefficient}, locally uniformly with two
source derivatives.  Consequently the exact-number and microcanonical
pressures exist and are obtained by the joint Legendre transform, not assumed
in advance.
\end{theorem}

\begin{proof}
Use Cauchy's formula in the activity and Fourier inversion in the constraints.
Move both contours through the real saddle.  The B2 normal analytic series
supplies a common complex neighborhood and a uniform cubic expansion.  The
activity direction is handled by the same fixed-size particle-number block,
and Theorem~\ref{thm:r21-b1-fourier} controls the remaining Fourier directions.
The local integral is the joint Gaussian determinant; every complementary
contour is exponentially smaller.  Uniform source derivatives follow by
marking the convergent B2 clusters.  Taking the logarithm and dividing by the
speed gives the claimed pressures.
\end{proof}

\begin{corollary}[Microcanonical transfer and Schur covariance]
\label{cor:r21-b1-transfer}
For every B2 path/contact source, conditioning on an admissible exact-number
microcanonical shell transfers the grand-canonical Laplace principle.  The
conditioned covariance is
\[
 \Sigma_{YY}-\Sigma_{YC}\Sigma_{CC}^{-1}\Sigma_{CY}
\]
on the quotient constraint space, and all source derivatives are evaluated at
the reoptimized saddle $\lambda(H,a)$.
\end{corollary}

\begin{proof}
Apply Theorem~\ref{thm:r21-b1-extraction} to numerator and denominator.  The
same local coefficient cancels, leaving the constrained pressure.  Differentiate
the saddle equations with respect to the source.  The implicit derivative of
$\lambda(H,a)$ produces exactly the displayed Schur complement.  Since
$\Sigma_{CC}$ is uniformly positive on $E$, the formula is well typed.
\end{proof}

\begin{theorem}[Round-twenty-one B1 closure]
\label{thm:r21-b1-main}
The positive exact-$N$ ensemble has a rank-correct global saddle, a fixed-size
block Fourier majorant covering velocity and position constraints, a precise
mixed lattice--continuous shell theorem, and a source-uniform coefficient
extraction from the B2 grand-canonical pressure.  No $C^4$ density is used for
$O(N)$ integrations by parts.
\end{theorem}

\begin{proof}
Combine Theorems~\ref{thm:r21-b1-saddle},
\ref{thm:r21-b1-fourier}, \ref{thm:r21-b1-coefficient}, and
\ref{thm:r21-b1-extraction}.
\end{proof}
